\documentclass[11pt]{article}

\usepackage[margin=1in]{geometry}
\usepackage{amsmath, amssymb}
\usepackage{graphicx}
\usepackage{booktabs}
\usepackage{hyperref}
\usepackage{xcolor}
\usepackage{enumitem}
\usepackage{caption}

\newcommand{\lcq}[1]{\textcolor{red}{#1}}

\title{AI-Assisted Workflow Optimisation\\for the HMM Transformer Research Project}
\author{C.~L.\ (with Claude Code)}
\date{\today}

\begin{document}
\maketitle
\tableofcontents
\bigskip

%----------------------------------------------------------------------
\section{Introduction}
%----------------------------------------------------------------------

This document records a set of workflow recommendations distilled from
accumulated interactions between the researcher and Claude Code over the
course of the HMM--transformer research project.  It is intended as a
living reference: a compact statement of conventions, guardrails, and
automation opportunities that have been identified as useful for this
particular combination of researcher background (physics), computing
environment (Linux cluster accessed via SSH from macOS), and project
structure (data generation $\to$ training $\to$ analysis pipeline).

The document is organised as follows.  Section~\ref{sec:env} describes
the computing environment and the role of the AI assistant within it.
Section~\ref{sec:coding} covers Python and PyTorch coding conventions
enforced by the assistant.  Section~\ref{sec:latex} covers \LaTeX{}
editing conventions.  Section~\ref{sec:cluster} covers cluster and
sweep submission practices.  Section~\ref{sec:general} states general
workflow preferences.  Section~\ref{sec:automation} discusses
automation opportunities opened up by access to a higher-capability
model tier (Claude Max / Opus 4.6).

%----------------------------------------------------------------------
\section{Environment and Role of the AI Assistant}
\label{sec:env}
%----------------------------------------------------------------------

The researcher works on a Linux cluster via SSH from a macOS laptop.
The following conventions apply.

\begin{itemize}[leftmargin=1.5em]
  \item \textbf{Command suggestions} for remote (cluster) tasks are given
        as Linux/Bash terminal commands.
  \item \textbf{Keyboard shortcuts} for local (macOS) tasks use Mac
        notation (\texttt{Cmd}, \texttt{Opt}, etc.).
  \item The assistant \textbf{never assumes a local desktop GUI
        environment}; all file editing, job submission, and inspection
        happen in the terminal or via SSH.
  \item The project lives at
        \texttt{\textasciitilde/workspace/hidden-markov-chain/};
        paths in scripts and configs should resolve from that root.
\end{itemize}

The AI assistant (Claude Code) acts as a collaborative coding partner
with three roles in this project:
\begin{enumerate}[leftmargin=1.5em]
  \item \textbf{Implementation} --- writing and running code, generating
        config files, submitting cluster jobs.
  \item \textbf{Explanation} --- translating ML conventions and
        mathematical results into language accessible to a physicist.
  \item \textbf{Documentation} --- producing \LaTeX{} reports,
        maintaining \texttt{CLAUDE.md} and \texttt{MEMORY.md} so that
        context persists across sessions.
\end{enumerate}

%----------------------------------------------------------------------
\section{Python and PyTorch Conventions}
\label{sec:coding}
%----------------------------------------------------------------------

\subsection{Device Consistency}

When writing code that mixes PyTorch tensors from different sources
(e.g.\ a tensor loaded from disk and one created in-flight during
training), device placement (CPU vs.\ CUDA) must be made explicit.

\begin{itemize}[leftmargin=1.5em]
  \item Use \texttt{.cpu()} before any NumPy comparison or
        \texttt{assert} that would fail silently if the tensor is on
        GPU.
  \item Use \texttt{.to(device)} consistently when passing tensors into
        model forward passes.
  \item Never rely on implicit broadcasting between a CPU and a CUDA
        tensor; PyTorch will raise a \texttt{RuntimeError}, but only at
        runtime, so the error may be hard to locate.
\end{itemize}

This convention was added explicitly because a physicist unfamiliar
with the GPU memory model may not think to check device placement when
a numerical result looks incorrect.

\subsection{Minimal Complexity Principle}

The assistant is instructed to prefer the simplest implementation that
satisfies the current requirement.  Concretely:
\begin{itemize}[leftmargin=1.5em]
  \item Do not add error handling for scenarios that cannot occur given
        the current code structure.
  \item Do not introduce helper abstractions for one-off operations.
  \item Do not add docstrings or type annotations to code that was not
        modified in the current task.
\end{itemize}
This prevents unnecessary growth of the codebase and keeps PRs focused.

%----------------------------------------------------------------------
\section{\LaTeX{} Editing Conventions}
\label{sec:latex}
%----------------------------------------------------------------------

When editing \texttt{.tex} files, the following checks are performed
automatically by the assistant.

\begin{enumerate}[leftmargin=1.5em]
  \item \textbf{Package dependencies}: any newly introduced symbol or
        environment must have its required package loaded in the
        preamble.  Common cases that have arisen in this project:
        \begin{itemize}
          \item \texttt{\textbackslash gtrsim}, \texttt{\textbackslash
                lesssim} $\to$ \texttt{amssymb}
          \item \texttt{align}, \texttt{gather} environments $\to$
                \texttt{amsmath}
          \item \texttt{tikzpicture} $\to$ \texttt{tikz}
          \item \texttt{tabularx}, \texttt{longtable} $\to$
                corresponding packages
        \end{itemize}
  \item \textbf{Test compilation}: after non-trivial edits, attempt
        \texttt{pdflatex} compilation to catch package or syntax errors
        before the changes are finalised.
  \item \textbf{Inline comments}: user comments of the form
        \verb|\lcq{...}| are addressed before the document is
        considered complete.  The comment markup is removed only upon
        explicit user approval.
\end{enumerate}

The figure path convention in this project is
\verb|\graphicspath{{../../figures/<experiment>/}}|, placing figures
in a top-level \texttt{figures/} directory organised by experiment
name, separate from the \texttt{notes/} source tree.

%----------------------------------------------------------------------
\section{Cluster, SLURM, and Sweep Submission}
\label{sec:cluster}
%----------------------------------------------------------------------

The cluster uses \texttt{addqueue} for job submission.  When generating
SLURM scripts or \texttt{wandb} sweep configs, the following checklist
is applied.

\begin{enumerate}[leftmargin=1.5em]
  \item \textbf{YAML numeric types}: values that look like numbers but
        should be treated as strings (e.g.\ experiment tags, run
        labels) must be quoted; otherwise the YAML parser silently
        converts them.
  \item \textbf{Variable substitutions}: placeholders such as
        \texttt{\$\{args\}} in sweep agent launch commands must be
        present and syntactically correct.
  \item \textbf{Environment loading}: module load commands and
        \texttt{conda activate} / \texttt{source .venv/bin/activate}
        lines must be included in every submission script, since
        cluster compute nodes do not inherit the login shell environment.
  \item \textbf{File path resolution}: all paths in scripts
        (\texttt{save\_dir}, \texttt{data\_dir}, config paths) are
        verified to resolve correctly from the cluster filesystem root,
        not from a local working directory.
\end{enumerate}

The main sweep tooling currently in use is \texttt{wandb sweep} with
\texttt{bash submit\_sweep.sh <sweep-id> [num\_agents]}.  All sweep
configs live under \texttt{config/sweeps/}.

%----------------------------------------------------------------------
\section{General Workflow Preferences}
\label{sec:general}
%----------------------------------------------------------------------

\begin{description}[leftmargin=1.5em, style=nextline]
  \item[Implement, don't only plan.]
        When asked to implement something, the assistant implements and
        runs it.  A plan document is produced only if the user
        explicitly requests it.  If a sweep or job submission is part
        of the task, it is submitted in the same session.

  \item[Ask for exact paths when searching fails.]
        If a first search attempt for a user-referenced file or
        function fails, the assistant asks the user for the exact file
        path rather than guessing at directory structures.  This avoids
        silent errors from editing the wrong file.

  \item[Flag non-standard ML practices.]
        Because the researcher's background is in physics, not machine
        learning, the assistant explicitly flags when a suggested
        approach deviates from standard ML practice, and briefly
        explains what is commonly done in the field.

  \item[Explain without coding by default.]
        When the user says ``explain'', the assistant provides a
        conceptual explanation only.  Code is written only if the user
        explicitly requests it in the same prompt.
\end{description}

%----------------------------------------------------------------------
\section{Automation Opportunities (Claude Max / Opus 4.6)}
\label{sec:automation}
%----------------------------------------------------------------------

Upgrading to Claude Max provides access to the Opus 4.6 model with a
200\,000-token context window and higher rate limits.  The following
automation opportunities are directly enabled or substantially improved
by this upgrade.

\subsection{Whole-Codebase Context}

With 200k tokens, the entire \texttt{src/} directory, all config
files, and the active \LaTeX{} report can be loaded simultaneously.
This eliminates the need to manually identify which files are relevant
to a given task, and allows the assistant to detect cross-file
inconsistencies (e.g.\ a config constraint violation before a training
run is submitted).

\subsection{End-to-End Experiment Automation}

A single assistant session can now cover a full experiment cycle:
\begin{enumerate}[leftmargin=1.5em]
  \item Edit \texttt{base\_config.yaml} for the new experiment.
  \item Generate data via \texttt{mpirun data\_generator.py}.
  \item Submit training sweep via \texttt{wandb sweep} and
        \texttt{submit\_sweep.sh}.
  \item Monitor W\&B run completion and pull results.
  \item Run analysis scripts and produce figures.
  \item Draft the \LaTeX{} section summarising the results.
\end{enumerate}
Steps 1--3 and 5--6 can be executed within one session without
context loss.

\subsection{Report Drafting and Revision}

The assistant can hold the full report \texttt{.tex} file in context
alongside the analysis outputs, allowing it to insert numerical
results, update figures, and address \verb|\lcq{}| comments in a
single pass.  This is particularly useful for the iterative
\texttt{time\_reversal/report.tex} workflow where multiple rounds of
edits are needed.

\subsection{Multi-Step Debugging Without Re-Explanation}

Longer context means the assistant retains the full stack trace,
config, and relevant source code across multiple back-and-forth turns
of debugging, without the user needing to re-paste context.  This
materially speeds up debugging cluster job failures and CUDA
environment issues.

\subsection{Parallel Sub-Agents (Agentic Mode)}

The Max plan enables launching multiple sub-agents in parallel.
Potential uses in this project:
\begin{itemize}[leftmargin=1.5em]
  \item Run \texttt{forward} and \texttt{reversed} training sweeps
        concurrently and compare results.
  \item Simultaneously search the codebase for two independent
        issues.
  \item Generate figures and draft the corresponding \LaTeX{}
        description in parallel.
\end{itemize}

\subsection{Persistent Memory and CLAUDE.md Maintenance}

With richer context, the assistant can more reliably detect when a new
convention or experimental result should be recorded in
\texttt{MEMORY.md} or \texttt{CLAUDE.md}, keeping the project
knowledge base up to date without explicit prompting.

%----------------------------------------------------------------------
\section{Summary and Next Steps}
\label{sec:summary}
%----------------------------------------------------------------------

Table~\ref{tab:recommendations} provides a concise summary of the
recommendations documented here.

\begin{table}[h]
\centering
\caption{AI workflow recommendations for the HMM transformer project.}
\label{tab:recommendations}
\begin{tabular}{lll}
\toprule
Category & Recommendation & Rationale \\
\midrule
Environment     & Use Linux commands for cluster tasks & SSH-only access \\
PyTorch         & Explicit \texttt{.to(device)} / \texttt{.cpu()} & Silent GPU/CPU mismatches \\
\LaTeX{}        & Check packages; compile after edits & Catch errors early \\
Cluster/Sweeps  & Verify YAML types, paths, env loads & Cluster env differs from login shell \\
Workflow        & Implement and run, not just plan & Reduces idle planning overhead \\
Workflow        & Ask for exact path if search fails & Avoids editing wrong files \\
Max plan        & Load whole codebase in context & Cross-file consistency checks \\
Max plan        & Parallel sub-agents for sweeps & Forward/reverse in one session \\
\bottomrule
\end{tabular}
\end{table}

Immediate next steps for adopting this workflow:
\begin{enumerate}[leftmargin=1.5em]
  \item Add a \texttt{Makefile} or shell script that runs the full
        data $\to$ train $\to$ analyse pipeline with a single command,
        so the assistant can trigger it without constructing the
        individual commands each time.
  \item Standardise figure generation scripts under
        \texttt{src/figures/} so the assistant knows where to look and
        what to run to regenerate any figure.
  \item Consider a lightweight \texttt{results\_registry.yaml} that
        records the W\&B run ID, config hash, and key metrics for each
        completed experiment, serving as a machine-readable index the
        assistant can query.
\end{enumerate}

\end{document}
